
% packages
\usepackage{mathtools}
\usepackage{xparse}
\usepackage{mathpartir}
\usepackage{etoolbox}
\usepackage{marvosym}
\usepackage{stmaryrd}
\usepackage{dsfont}
\usepackage{minted}
\usepackage{hyperref}
\usepackage{mathtools}
\usepackage{varwidth}
\usepackage{bussproofs}
\usepackage{framed}
\usepackage{caption}
\usepackage{amssymb}


% utils
\newcommand\ifempty[3]{\expandafter\ifstrempty{#1}{#2}{#3}}
\newcommand\ifnempty[2]{\ifempty{#1}{}{#2}}

\newcommand\ifthenabove[1]{%
    \ifnempty{#1}{^{#1}}%
}

\makeatletter
\newcommand\mpar@math[3]{\ensuremath{\left(#1\right)\ifnempty{#2}{_{#2}}\ifnempty{#3}{^{#3}}}}
\newcommand\mpar@text[3]{\ensuremath{(#1)\ifnempty{#2}{_{#2}}\ifnempty{#3}{^{#3}}}}
\newcommand\mpar@belowabove[3]{\mathchoice{\mpar@math{#1}{#2}{#3}}{\mpar@text{#1}{#2}{#3}}{\mpar@text{#1}{#2}{#3}}{\mpar@text{#1}{#2}{#3}}}

\def\mpar@grabbelow#1#2#3_#4{%
  \mpar@grab{\ifnempty{#1}{#1, }#4}{#2}{#3}%
}
\def\mpar@grabbove#1#2#3^#4{%
  \mpar@grab{#1}{\ifnempty{#2}{#2, }#4}{#3}%
}
\newcommand\mpar@grab[3]{%
  \@ifnextchar_{\mpar@grabbelow{#1}{#2}{#3}}%
  {\@ifnextchar^{\mpar@grababove{#1}{#2}{#3}}%
  {\mpar@belowabove{#3}{#1}{#2}}}%
}

\newcommand\mpar{\mpar@grab{}{}}
\makeatother

\newcommand\funp{\mpar}
\newcommand\funu[1]{\ifnempty{#1}{\funp{#1}}}
\newcommand\fund[2]{\ifnempty{#1#2}{\funp{\sep{#1}{#2}}}}
\newcommand\funt[3]{\ifnempty{#1#2#3}{\funp{\sept{#1}{#2}{#3}}}}
\newcommand\funq[4]{\ifnempty{#1#2#3#4}{\funp{\sepq{#1}{#2}{#3}{#4}}}}

\newcommand\sep[2]{#1\ifnempty{#1}{\ifnempty{#2}{,}}#2}
\newcommand\sept[3]{\sep{#1}{#2}\ifnempty{#1#2}{\ifnempty{#3}{,}}#3}
\newcommand\sepq[4]{\sept{#1}{#2}{#3}\ifnempty{#1#2#3}{\ifnempty{#4}{,}}#4}

\newcommand{\eqdef}{\triangleq}

% comments

\usepackage{todonotes}
\newcommand\todoi[1]{\todo[inline,caption={}]{#1}}
\newcommand{\as}[1]{\todo[color=blue!30,caption={},size=\tiny]{#1}}
\newcommand{\asi}[1]{\todo[color=blue!30,inline]{#1}}
\renewcommand{\ln}[1]{\todo[color=magenta!50,caption={},size=\tiny]{#1}}
\newcommand{\lni}[1]{\todo[color=magenta!50,inline]{#1}}

% lists and sets
\newcommand\cons[2]{#1; #2}
\newcommand\emptylist{[]}
\newcommand\dlist[1]{\left[#1\right]}
\newcommand\singlelist[1]{\dlist{#1}}
\newcommand\lconcat{+\!\!+}
\newcommand\lcons{{::}}

\newcommand\mset[2][]{%
    \ifempty{#2#1}%
	    {\emptyset}%
	    {\left\{%
		    %\begin{gathered}{c}%
            #1%
            \ifnempty{#1}{\middle|}%
		    #2%
		    %\end{gathered}%
	    \right\}}%
}

% tuples

% \newcommand\myvecbase[3][]{\overrightarrow{\ifnempty{#1}{#1\funu}{#2_{#3}}}}

% Given:
% #1 (facultative): an identifier, meant to mean a function (if absent, it will just not be written),
% #2: a function taking a number.
% #3: the maximum number.
% Generates #1(#2{1})..#1(#2{#3})
\newcommand\myvecbase[3][]{\ifnempty{#1}{#1\funu}{#2{1}}..\ifnempty{#1}{#1\funu}{#2{#3}}}

% The p and s variants respectively add parentheses and braces.
\newcommand\myvecm{\myvecbase}
\newcommand\myvecmp[3][]{\mpar{\myvecm[#1]{#2}{#3}}}
\newcommand\myvecms[3][]{\mset{\myvecm[#1]{#2}{#3}}}

% Same, but the #2 function takes a facultative argument (possibly followed
% by a third mandatory one, which will be instanciated empty).
\newcommand\mandatorytofacultative[2]{#1[#2]{}}
\newcommand\myvec[2][]{\myvecbase[#1]{\mandatorytofacultative{#2}}}
\newcommand\myvecp[3][]{\mpar{\myvec[#1]{#2}{#3}}}
\newcommand\myvecs[3][]{\mset{\myvec[#1]{#2}{#3}}}

% Same, but #2 is now an identifier, which will be subscripted the index.
\newcommand\mandatorytoindex[2]{#1\ifnempty{#2}{_{#2}}}
\newcommand\myveci[2][]{\myvecbase[#1]{\mandatorytoindex{#2}}}
\newcommand\myvecip[3][]{\mpar{\myveci[#1]{#2}{#3}}}
\newcommand\myvecis[3][]{\mset{\myveci[#1]{#2}{#3}}}

% Given:
% #3: an equivalent of \mapsto
% #1: the “from” part of the mapping.
% #2: the “from” internal command, taking facultatively an index.
% #4: the “to” part of the mapping.
% #5: the “to” internal command, taking facultatively an index.
% #6: the maximum number.
% Generates a sequence of mappings #2[1] \mapsto #5[1]..#2[#6] \mapsto #5[#6],
% or #1(#2[1]) \mapsto #1(#5[1])..#1(#2[#6]) \mapsto #1(#5[#6]),
\makeatletter
\newcommand\myvecmap@fun[6]{\ifnempty{#1}{#1\funu}{#2[#6]{}} #3 \ifnempty{#4}{#4\funu}{#5[#6]{}}}
\newcommand\myvecmap[6]{%
    \myvecbase{\myvecmap@fun{#1}{#2}{#3}{#4}{#5}}{#6}}
%   \ifnempty{#1}{#1\funu}{#2[1]{}} #3 \ifnempty{#4}{#4\funu}{#5[1]{}}%
%   ..\ifnempty{#1}{#1\funu}{#2[#6]{}} #3 \ifnempty{#4}{#4\funu}{#5[#6]}}
\makeatother

% misc
\newcommand\dom{\mathit{dom}\funu}
\newcommand\true{\ensuremath{\top}}

% variables
\newcommand{\tvar}[1][]{x_{t\ifnempty{#1}{_{#1}}}}
\newcommand{\fvar}[1][]{x_{f\ifnempty{#1}{_{#1}}}}
\DeclareDocumentCommand{\svar}{ m O{} }{#1{\ifnempty{#2}{_{#2}}}}
\newcommand{\ivar}[1][]{x_{\sigma\ifnempty{#1}{_#1}}}
\newcommand{\ovar}{x_o}

% Sorts and Terms
\newcommand\Sort[1][\typenv]{\ensuremath{\mathit{Sort}_{#1}\funu}}
\newcommand\sort[1][]{\ensuremath{s\ifnempty{#1}{_{#1}}}}
\newcommand\sortp[1][]{\ensuremath{s'\ifnempty{#1}{_{#1}}}}
\newcommand\Tvar{\ensuremath{\mathit{Tvar}\funu}}
\newcommand\sig{\ensuremath{\mathit{sig}\funu}}
\newcommand{\termenv}{E}
\newcommand{\term}[1][]{t\ifnempty{#1}{_{#1}}}
\newcommand{\termp}[1][]{t'\ifnempty{#1}{_{#1}}}
\newcommand{\progpoint}{\texttt{pp}}
\newcommand\ppe{\epsilon}
\newcommand\ppc[3][]{\ifnempty{#2}{#2}{\cdot}\ifnempty{#3}{\ifempty{#1}{#3}{#3[#1]}}}
\newcommand\cppv[3][]{\ifnempty{#2}{#2}{\rule[0.5ex]{2pt}{0.5pt}}\ifnempty{#3}{\ifempty{#1}{#3}{#3[#1]}}}
\newcommand\syntaxcontraint[1]{\left[#1\right]}
\newcommand\subterm[2]{#1@#2}
\newcommand{\ppenv}{\mathit{HPP}}
\newcommand{\match}[4]{#1 = #3\funu{#2}:#4}

% Rules
\newcommand{\bvar}{V}

\newcommand\rulen[1]{\textit{#1}}
\newcommand\Rule[4]{\ensuremath{\rulen{#1}{\mpar{#2}} \coloneqq #3_{#4}}}
\newcommand\Rulea[4]{\ensuremath{\rulen{#1}{\mpar{#2}} & \coloneqq #3_{#4}}} % Only to be used in an {align*} environment.
\newcommand{\Rules}{\ensuremath{\mathit{Rules}}}


\newcommand\derivH{H}
\newcommand\letin[2]{\operatorname{let}#1=#2\operatorname{in}}
\newcommand\skexist[2]{\operatorname{let}#1:#2\operatorname{in}}
\newcommand\deriv[2]{{\texttt{#1}}~{#2}}
\newcommand\branches[1]{\mpar{#1}}
\newcommand\branchesV{\branches}
\newcommand{\eqp}{\mathrel{?\!\vartriangleright}}
\newcommand\filter[3]{#1\funu{#2}\ifnempty{#3}{\eqp #3}}

\newcommand{\bcons}[2]{#1 \mathrel{\Arrowvert} #2}
\newcommand{\bnil}{{\bullet}}

\newcommand{\skel}[1][]{S{\ifnempty{#1}{_{#1}}}}
\newcommand{\bone}[1][]{B{\ifnempty{#1}{_{#1}}}}

\newcommand\judgment[3]{\ifnempty{#1}{#1\ifnempty{#2}{, }}#2 \Downarrow #3}

% While Rules
%% Sorts
\newcommand\litS{\ensuremath{\mathit{lit}}}
\newcommand\identS{\ensuremath{\mathit{ident}}}
\newcommand\exprS{\ensuremath{\mathit{expr}}}
\newcommand\statS{\ensuremath{\mathit{stat}}}
\newcommand\stackS{\ensuremath{\mathit{stack}}}
\newcommand\valS{\ensuremath{\mathit{val}}}

%% Syntax
\newcommand\jsid[2][]{{\textup{\texttt{\small #2}}\ifthenabove{#1}}}
\newcommand\Var{\ensuremath{\mathit{Var}}}
\newcommand\constn{\mathit{const}} 
\newcommand\const[1]{\constn\funu{#1}}
\newcommand\varn{\mathit{var}} 
\newcommand\var[1]{\varn\funu{#1}}
\newcommand\access[1]{!#1}
\newcommand{\erefn}{\ensuremath{\mathit{ref}}}
\newcommand\eref[1]{\erefn\funu{#1}}
\newcommand{\ein}{\ensuremath{\mathit{in}}}
\newcommand{\eoutn}{\ensuremath{\mathit{out}}}
\newcommand{\eout}[1]{\eoutn\mpar{#1}}
\newcommand\sskip{\mathit{skip}}
\newcommand\sthrow{\mathit{throw}}
\newcommand\iname{\mathit{if}}
\newcommand\wname{\mathit{while}}
\newcommand\asnescape[2]{#1\ {:=}\ #2}
\newcommand{\heapwrite}[2]{#1 \leftarrow #2}
\newcommand\seq[2]{{#1};\ #2}
\newcommand\stry[2]{\mathit{try} \, #1 \, \mathit{catch} \, #2}
\newcommand\sif[3]{\iname\ifnempty{#1#2#3}{\,#1\,#2\ifnempty{#3}{\,#3}}}
\newcommand\while[2]{\wname\ifnempty{#1#2}{\,#1\,#2}}
\newcommand{\progvar}{\jsid{Var}}
\newcommand{\proglit}{\jsid{Lit}}

%% Interpretations
\newcommand{\interpIn}{\Sigma}
\newcommand{\interpOut}{O}
\newcommand\vdefined[1]{{#1\downarrow}}
\DeclareDocumentCommand{\interpDef}{ O{\interpIn} O{I} m }{\ensuremath{\vdefined{\left\llbracket #1, #3 \right\rrbracket^{#2}}}}
\DeclareDocumentCommand{\interp}{ O{\interpIn} O{I} m O{\interpOut}}{\ensuremath{%
    \left\llbracket #3 \right\rrbracket\ifnempty{#2}{^{#2}}\ifnempty{#1}{\mpar{#1}}%
    \ifnempty{#4}{ \Downarrow #4}}}

\newcommand{\interpbIn}[1][\bvar]{\interpIn_{#1}}
\newcommand{\interpbOut}[1][\bvar]{\interpOut_{#1}}
\DeclareDocumentCommand{\interpbDef}{ O{\interpIn} O{I} m O{\bvar}
}{\ensuremath{\vdefined{\left\llbracket #3
\right\rrbracket^{#2}_{#4}\mpar{#1}}}}
\DeclareDocumentCommand{\interpb}{ O{\interpIn} O{I} m O{\interpOut}
  O{\bvar}}{\ensuremath{\left\llbracket #3 \right\rrbracket\ifnempty{#2}{^{#2}}\ifnempty{#5}{_{#5}}\mpar{#1}
    \Downarrow #4}}

\DeclareDocumentCommand{\merge}{ O{I} m m O{\interpIn'} O{\bvar}
}{\ensuremath{\interpb[#2][#1]{\bigoplus_{#3}}[#4][#5]}}
\newcommand{\outfun}{\mathcal{O}}
\newcommand{\outfuna}{\abs{\mathcal{O}}}

\newcommand{\csem}{\ensuremath{\Downarrow}}
\newcommand{\asem}{\ensuremath{\abs{\csem}}}
\newcommand{\ssem}{\ensuremath{\abs{\csem_{s}}}}

\DeclareDocumentCommand{\typvar}{ m O{} }{\ensuremath{\tau\ifnempty{#1}{#1}\ifnempty{#2}{_{#2}}}}
\newcommand\typenv{\Gamma}
\newcommand\extsymbol{ + }
\newcommand\extenv[3]{#1\extsymbol#2\leftarrow#3}
\newcommand\extenvn[4]{#1\extsymbol\mpar{#2_1 \mapsto#3_1..#2_{#4} \mapsto#3_{#4}}}
\newcommand\envrestr[2][\bvar]{\left.#2\right|_{#1}} 
\newcommand\typin{\mathit{in}\funu}
\newcommand\typout{\mathit{out}\funu}
\newcommand{\typsig}{\mathit{fsort}\funu}
\newcommand{\typinterp}[4][V]{#2 \vdash_{#1} #3: #4}
\newcommand{\dfvarset}{\ensuremath{\mathcal{D}}}
\DeclareDocumentCommand{\wfinterp}{ O{\typenv,\dfvarset} m
  O{\mpar{\typenv',\dfvarset'}} O{} }{\interpb[#1][\texttt{wf}]{#2}[#3][#4]}

\newcommand\valT{\mathit{val}}
\newcommand\intT{\mathit{int}}
\newcommand\boolT{\mathit{bool}}
\newcommand\locT{\mathit{loc}}
\newcommand\stateT{\mathit{store}}
\newcommand\heapT{\mathit{heap}}
\newcommand\IOstateT{\mathit{state}}
\newcommand\valIOstateT{\mathit{valState}}
\newcommand\excIOstateT{\mathit{excState}}
\newcommand\inT{\mathit{in}}
\newcommand\outT{\mathit{out}}
\newcommand\unit{()}
\newcommand\resundef{\bot}

\newcommand{\concrenv}{\Sigma}
\newcommand{\quadset}{Q}

\newcommand\Some{\texttt{Some}\funu}
\newcommand\None{\texttt{None}}

% Concrete interpretation of a filter, expressed as a relation.
\DeclareDocumentCommand{\concrinterp}{ O{\concrenv,\quadset} m
  O{\mpar{\concrenv',\quadset}} O{} }{ \interpb[#1][]{#2}[#3][#4] }

\DeclareDocumentCommand{\filterinterp}{ O{} m m O{\Downarrow} m }{\ensuremath{\left\llbracket #2
\right\rrbracket^{#1}\ifnempty{#3}{\funu{#3}\ifnempty{#5}{#4}}#5}}

\newcommand{\concrfilter}[3]{\filterinterp[]{#1}{#2}{#3}}
\newcommand{\concrfilterDef}[2]{\vdefined{\filterinterp{#1}{#2}[]{}}}
\newcommand{\notconcrfilter}[2]{\filterinterp{#1}{#2}[\not\Downarrow]{}}
\newcommand{\cinterpf}[1][]{\immediateconsequence\ifnempty{#1}{^{#1}}\funu}

\newcommand{\cval}[1][]{v\ifnempty{#1}{_{#1}}}
\newcommand{\cvalp}[1][]{v'\ifnempty{#1}{_{#1}}}

\newcommand{\aterm}[1][]{\term^{\#}\ifnempty{#1}{_{#1}}}
\newcommand{\atermp}[1][]{\term^{\#\prime}\ifnempty{#1}{_{#1}}}

\newcommand{\aval}[1][]{v^{\#}\ifnempty{#1}{_{#1}}}
\newcommand{\avalp}[1][]{v\ifnempty{#1}{_{#1}}^{\#\prime}}
\newcommand{\avalpp}[1][]{v\ifnempty{#1}{_{#1}}^{\#\prime\prime}}

\newcommand{\cstate}[1][]{\sigma\ifnempty{#1}{_{#1}}}
\newcommand{\astate}[1][]{\sigma^{\#}\ifnempty{#1}{_{#1}}}
\newcommand{\astatep}[1][]{\sigma^{\#\prime}\ifnempty{#1}{_{#1}}}
\newcommand{\astatepp}[1][]{\sigma^{\#\prime\prime}\ifnempty{#1}{_{#1}}}

\newcommand{\abstrint}[3]{\interp[#1][\#]{#2}[#3]}

% Abstract interpretation of a skeleton or bone (mandatory argument).
% The facultative arguments describe the input and output states
% The last facultative argument is an index of the abstract interpretation.
% Note about the \absflag argument that indicates whether we are missing this branch.
\DeclareDocumentCommand{\abstrinterp}{ O{\absflag} O{\abstrenv,\abstripleset}  m
  O{\absflag} O{\abstrenvp,\abstripleset} O{} }{ \interpb[#1,#2][\#]{#3}[\mpar{#4,#5}][#6] }

\newcommand{\concr}{\gamma}

\newcommand{\abstrintDef}[3][\bvar]{\interpbDef[#2][\#]{#3}[#1]}
\newcommand{\abstrenv}[1][]{\Sigma^\#\ifnempty{#1}{_{#1}}}
\newcommand{\abstrenvp}[1][]{\Sigma^{\#\prime}\ifnempty{#1}{_{#1}}}
\newcommand{\abstrenvpp}[1][]{\Sigma^{\#\prime\prime}\ifnempty{#1}{_{#1}}}
\newcommand\abstripleset[1][]{{T^\#\ifnempty{#1}{_{#1}}}} % An abstract set of triples.
\newcommand{\abstrfilter}[3]{\filterinterp[\#]{#1}{#2}[=]{#3}}
\newcommand{\abstrfilterDef}[2]{\vdefined{\filterinterp[\#]{#1}{#2}[]{}}}

% Given:
% #2: The argument of the immediate consequence (a set of triples)
% #1 (facultative): an exponent, to apply the immediate consequence several times.
% Builds the abstract immediate consequence.
\newcommand\ainterpf[1][]{{\immediateconsequence^\#}\ifnempty{#1}{^{#1}}\funu}

\newcommand{\stSplit}{\mathit{Sp}\funu}
\newcommand\sinterpf[1]{\stSplit{\ainterpf{\stSplit{#1}}}}

\newcommand\Kcombinator[2]{#1}

\makeatletter
\def\abot@complete#1[#2]{\bot{\ifempty{#2}{{\ifnempty{#1}{_{#1}}}}{_{#1[#2]}}}}
\newcommand\abot[1]{\@ifnextchar[{\abot@complete{#1}}{\abot@complete{#1}[]}}
\makeatother

%% Consistency
\newcommand\OKst[1][]{\mathit{OKst}\ifnempty{#1}{_{#1}}\fund}
\newcommand\OKout[1][]{\mathit{OKout}\ifnempty{#1}{_{#1}}\fund}

%% Equation Generation

\newcommand{\ppgen}[1][t_0]{\ppenv_{#1}\funt}
\newcommand{\eqset}{\mathcal{C}}
\newcommand{\eqenv}{D_e}
\newcommand{\eqdfvar}[1][N]{\mathcal{D}_{\rulen{#1}}\funu}
\newcommand{\eqgen}[3][\rulen{N},\progpoint,\eqset]{\interp[#1][c]{#2}[#3]}
\newcommand{\eqgenb}[3][\rulen{N},\progpoint,\eqset]{\interpb[#1][c]{#2}[#3]}
\newcommand{\eqfilter}[3]{\filterinterp[c]{#1}{#2}[=]{#3}}
\newcommand\Gen{\mathit{Gen}\funu}
\newcommand\Gent{\mathit{PP}\funu}
\newcommand{\bhist}{\texttt{bb}}
\newcommand\bbe{\epsilon}
\newcommand\bbc[2]{\ifnempty{#1}{#1}.\ifnempty{#2}{#2}}

\newcommand\svareq[3]{\svar{#1}[#3] \leftarrow {#2}\ifnempty{#3}{_{#3}}}
\newcommand\fvareq[3]{\fvar[#1#3]{} \leftarrow {#2}\ifnempty{#3}{_{#3}}}
\newcommand\tvareq[3]{\tvar[#1#3]{} \leftarrow {#2}\ifnempty{#3}{_{#3}}}

\newcommand\Val{\ensuremath{\mathit{Val}}}
\newcommand\val{\ensuremath{v}}

\makeatletter % There may be a simpler way to do this.
% See usage below.
\newcommand\@absstarnothing[1]{\ensuremath{\mathrel{#1^{\#}}}}
\newcommand\@absnostarnothing[1]{\ensuremath{{#1^{\#}}}}
\def\@absstarbrack#1[#2]{\ensuremath{\mathrel{#1^{\#}\left[#2\right]}}}
\def\@absnostarbrack#1[#2]{\ensuremath{{#1^{\#}\left[#2\right]}}}
\newcommand\@absstarnounder[1]{%
    \@ifnextchar[{\@absstarbrack{#1}}{\@absstarnothing{#1}}}
\newcommand\@absnostarnounder[1]{%
    \@ifnextchar[{\@absnostarbrack{#1}}{\@absnostarnothing{#1}}}
\def\@absstarunder#1_#2{\ensuremath{\mathrel{#1_{#2}^{\#}}}}
\def\@absnostarunder#1_#2{\ensuremath{{#1_{#2}^{\#}}}}
\newcommand\@absstarchoice[1]{%
    \@ifnextchar_{\@absstarunder{#1}}{\@absstarnounder{#1}}}
\newcommand\@absnostarchoice[1]{%
    \@ifnextchar_{\@absnostarunder{#1}}{\@absnostarnounder{#1}}}
\newcommand\abs{\@ifstar\@absstarchoice\@absnostarchoice}
% Usages:
%   \abs{value}
%   \abs{value}_index
%   \abs{value}[arg]
%   \abs*{relation}
%   \abs*{relation}_index
%   \abs*{relation}[arg]
\makeatother
\newcommand\bigabs[1]{\mathrel{\abs*{{#1}}}}

% Was \newcommand\absint[1]{\texttt{A}\funu{#1}} %
\newcommand\absint[1]{\llbracket #1 \rrbracket^{\#}}
%  need to change the concrete version accordingly?

\newcommand*\mathsc[1]{\mathrm{\textsc{#1}}}


\newcommand\itvl[2]{[ #1 , #2]}

\newcommand\ptrue{\textit{true}}
\newcommand\pfalse{\textit{false}}
\newcommand\btrue{\ptrue}%{\texttt{tt}}
\newcommand\bfalse{\pfalse}%{\texttt{ff}}

\newcommand\eread[2]{#1\funu{#2}}
\newcommand\ewrite[3]{#2\left[#1 \leftarrow #3\right]}

\newcommand\NSRname{\mathit{NS}}
\newcommand\NSR[1]{\NSRname\ifnempty{#1}{_{\mathsc{#1}}}\funu} % Non-structural rule.
\newcommand\immediateconsequence{\mathcal{H}}
\newcommand\immediateconsequenceglue[1][]{\abs{G\ifnempty{#1}{_{#1}}}}

\newcommand\evalr{\ensuremath{\Downarrow}}
\newcommand\cinterp[3][\evalr]{\ensuremath{\left\llbracket #2 \right\rrbracket_{#1}\funu{#3}}}
\newcommand\ainterp[3][\abs*\evalr]{\ensuremath{\abs{\left\llbracket #2 \right\rrbracket}_{#1}\funu{#3}}}

\newcommand\extsig[3]{\ifnempty{#1}{#1\extsymbol }#2 \mapsto #3}
\newcommand\extsign[3]{#1\extsymbol \overrightarrow{#2 \mapsto #3}}

% Given
% #1: an environment (which can be empty)
% #3 and #5: a command taking facultatively a number, respectively representing the key and the value.
% #2 and #4 (facultative): a symbol representing a function that will be respectively applied to the arguments #3 and #5.
% #3 (facultative): the number of terms to be added.
% Builds the extended environment.
\DeclareDocumentCommand{\extsigvec}{ m O{} m O{} m O{n} }{\ifnempty{#1}{#1\extsymbol}\myvecmap{#2}{#3}{\mapsto}{#4}{#5}{#6}}

\newcommand\outst[2]{\mathit{out}\funu{#1,#2}}

\newcommand{\While}{\textsc{While}}
\newcommand\wtrue{\texttt{true}}
\newcommand\wfalse{\texttt{false}}
\newcommand\wint{\texttt{n}}

\newcommand{\halt}{\bot}
\newcommand{\cont}{\top}
\newcommand{\absflag}{f}

\newcommand\K{\ensuremath{\mathbb{K}}}



%%% Local Variables:
%%% mode: latex
%%% TeX-master: "new_necro"
%%% End:
